\begin{abstract}
A procedure can be correct in the setting where it was learned and inappropriate in the setting where it is reused. We study this distinction in coding agents by supplying source-valid procedural memory while changing a security-relevant assumption in the target task. Thirteen synthetic task families cross four memory conditions with two repetitions in each of six model--agent configurations, yielding \RecordedRuns{} recorded runs and \ValidRuns{} technically valid outcomes. The primary endpoint is functional completion that fails a bounded security witness. Within-configuration comparisons do not establish a uniform harmful-memory effect. Correct-memory contrasts range from $\PrimaryURangeLow$ to $\PrimaryURangeHigh$ percentage points; lower failure rates can accompany reduced functionality. A generic applicability reminder reduces the endpoint descriptively in the three Codex configurations but has no consistent benefit in MiniSWE. Selected transcript--implementation--outcome comparisons show both procedure-consistent failures and successful adaptation, including cases where correct memory itself makes a changed assumption salient. These findings concern supplied memory in shifted synthetic contexts, not an end-to-end retrieval system. They show why memory transfer must be assessed jointly with functionality, witness scope, and the agent's execution interface.
\end{abstract}
